\subsection{Tests, Likelihood Ratios, and Error Control}

Testing formalizes decisions between hypotheses.  The central quantities are the probability of false rejection under the null and the probability of rejection under alternatives.

\begin{definition} \label{def:hypotheses-tests-size-power} \hlt{Hypotheses, Tests, Size, and Power}\\
A \hlt{null hypothesis} is a subset $\Theta_0\subseteq\Theta$, and an \hlt{alternative} is a subset $\Theta_1\subseteq\Theta$ disjoint from $\Theta_0$.  A randomized \hlt{test} is a measurable function $\varphi:\mathcal{X}\to[0,1]$, where $\varphi(x)$ is the rejection probability.  Its \hlt{size} and \hlt{power function} are
\[
\alpha(\varphi)=\sup_{\theta\in\Theta_0}\mathbb{E}_\theta[\varphi(X)],
\qquad
\beta_\varphi(\theta)=\mathbb{E}_\theta[\varphi(X)].
\]
\end{definition}

\begin{theorem} \label{thm:neyman-pearson-lemma} \hlt{Neyman--Pearson Lemma}\\
For testing a simple null $P_0$ against a simple alternative $P_1$, a most powerful level-$\alpha$ test rejects for large likelihood ratio
\[
\Lambda(X)=\frac{dP_1}{dP_0}(X),
\]
with randomization on the boundary if needed.
\begin{proof}
For any constant $c\geq0$ and any competing level-$\alpha$ test $\psi$, compare $\varphi-\psi$ on the regions $\Lambda>c$ and $\Lambda<c$.  The sign of $(\varphi-\psi)(\Lambda-c)$ is nonnegative for the likelihood-ratio rule, so
\[
\mathbb{E}_{1}[\varphi-\psi]
=\mathbb{E}_{0}[\Lambda(\varphi-\psi)]
\geq c\,\mathbb{E}_{0}[\varphi-\psi]\geq0
\]
after choosing the boundary randomization so that $\mathbb{E}_0\varphi=\alpha$.
\end{proof}
\end{theorem}

\begin{definition} \label{def:generalized-likelihood-ratio-test} \hlt{Generalized Likelihood-Ratio Test}\\
For composite hypotheses, the \hlt{generalized likelihood ratio} is
\[
\lambda(x)=
\frac{\sup_{\theta\in\Theta_0}L_x(\theta)}
{\sup_{\theta\in\Theta}L_x(\theta)}.
\]
Small values of $\lambda(x)$ are evidence against $\Theta_0$.
\end{definition}

\begin{definition} \label{def:p-value-and-validity} \hlt{$p$-Value}\\
A \hlt{$p$-value} is a statistic $p(X)\in[0,1]$ such that, for every $\theta\in\Theta_0$ and every $u\in[0,1]$,
\[
\mathbb{P}_\theta(p(X)\leq u)\leq u.
\]
Rejecting when $p(X)\leq\alpha$ therefore gives a level-$\alpha$ test.
\end{definition}

\begin{remark} \label{rem:error-control-families} \hlt{Error Rates}\\
Single-test size controls one false rejection probability.  In multiple testing, common targets include the family-wise error rate, the false discovery rate, and directional or selective error rates.  Each requires a statement about the family of decisions being protected.
\end{remark}
